\begin{solapartat}
\punts{0,65}

\noindent\begin{minipage}[c]{0.66\linewidth}
\figura[1]{sol_fig1.png}
\end{minipage}\hfill
\begin{minipage}[c]{0.3\linewidth}
\figura[1]{sol_fig2.png}
\end{minipage}

Les línies de camp magnètic són circumferències en un pla perpendicular al fil i centrades amb el fil, per tant:
\begin{itemize}
  \item cal dibuixar els camps perpendiculars als fils i en la direcció indicada en el dibuix,
  \item el sentit del camp es pot determinar segons la regle de la ma dreta.
\end{itemize}

\punts{0,6} El camp magnètic creat pel fil és perpendicular al fil. Com els fils són perpendiculars entre ells, el camp magnètic creat pel fil $A$, $\vec{B}_A$, és paral·lel al fil $B$ i viceversa. Llavors, el corrent que ha de circular pel fil $B$ ha de ser zero, d'aquesta manera només circularà corrent pel fil $A$, i el camp total serà el generat pel fil $A$, $\vec{B}_A$, que és paral·lel al fil $B$.
\end{solapartat}

\begin{solapartat}
% DUBTE: a la pauta diu «I com B_A i B_A són perpendiculars», per «B_A i B_B».
\punts{0,75} Si el camp magnètic total fa un angle de 30° respecte el fil $B$, llavors el camp magnètic total fa un angle de 30° respecte $\vec{B}_A$. I com $\vec{B}_A$ i $\vec{B}_A$ són perpendiculars entre ells:
\[ \frac{B_B}{B_A} = \tan(30^\circ) \]
\figura[0.15]{sol_fig3.png}

\punts{0,4} La intensitat de camp magnètic creada per un fil molt llarg en un punt és proporcional a la intensitat que hi circula i depèn de la distància d'aquest punt al fil. Com la distància és la mateixa pels dos fils, llavors el quocient d'intensitats de camp magnètic és igual al quocient de corrents:
\[ \frac{I_B}{I_A} = \frac{B_B}{B_A} = \tan(30^\circ) \]

\punts{0,1} $I_B = I_A\tan(30^\circ) = 5,77\ \mathrm{A}$
\end{solapartat}
